\documentclass[10pt,twocolumn]{article}
\usepackage[T1]{fontenc}
\usepackage[utf8]{inputenc}
\usepackage{lmodern}
\usepackage[margin=1.8cm,top=2.4cm,bottom=2cm,columnsep=0.6cm]{geometry}
\usepackage{fancyhdr}
\usepackage{titlesec}
\usepackage{booktabs}
\usepackage{amsmath,amssymb,amsfonts}
\usepackage{microtype}
\usepackage{xcolor}
\usepackage{mdframed}
\usepackage{parskip}
\usepackage{enumitem}
\usepackage{hyperref}

\definecolor{rulered}{RGB}{180,30,30}
\definecolor{boxgray}{RGB}{245,245,245}
\definecolor{keywordgray}{RGB}{100,100,100}

\pagestyle{fancy}
\fancyhf{}
\fancyhead[L]{\textsc{\small Claude Mag}}
\fancyhead[C]{\small\textit{Machine Learning Research Digest}}
\fancyhead[R]{\small 2026-08-10}
\fancyfoot[C]{\small\thepage}
\renewcommand{\headrulewidth}{0.8pt}

\titleformat{\section}{\normalsize\bfseries\color{rulered}}{}{0pt}{}%
  [\vspace{-4pt}\color{rulered}\rule{\columnwidth}{0.4pt}]
\titlespacing{\section}{0pt}{8pt}{4pt}

\hypersetup{colorlinks=true,urlcolor=rulered,linkcolor=black}

\begin{document}
\setlength{\parindent}{0pt}
\setlength{\parskip}{4pt}

{\small\color{keywordgray}\textbf{Keywords:}
  state space models \textbullet{}
  long-sequence modeling \textbullet{}
  hierarchical memory \textbullet{}
  recurrent neural networks}\par\vspace{4pt}

{\LARGE\bfseries\raggedright
  Mamba with Hierarchical Memory: Solving the
  Representation Bottleneck in Long Sequence Modeling}\par\vspace{4pt}

{\large\itshape\raggedright
  HMM Adds Paragraph-Level Semantics and Persistent Long-Term Memory
  to Mamba's Fixed Recurrent State, Enabling Near-Linear Long-Context
  Reasoning}\par\vspace{6pt}

{\small\textbf{By} Qinwen Wang, Jieping Luo, Aoxiang Qin, Ruoyu Zhao,
  Jianxiong Tang, Wei Zhang, Zhichao Lu,
  Luziwei Leng}\par\vspace{2pt}

{\small\color{keywordgray}arXiv:2608.02347 \textbullet{} August 2026}
\par\vspace{2pt}\noindent{\color{rulered}\rule{\columnwidth}{0.6pt}}\vspace{6pt}

State space models like Mamba offer linear-time sequence processing, but their
fixed-capacity recurrent state is a fundamental bottleneck on long sequences: old
context is progressively overwritten as the sequence grows. Hierarchical Memory Mamba
(HMM) resolves this by adding two lightweight memory tiers above the base recurrent
state---paragraph-level semantic extraction and persistent long-term memory retrieval---
enabling near-linear complexity long-range reasoning on sequences where Mamba's
performance degrades.

\begin{mdframed}[backgroundcolor=boxgray,linecolor=rulered,linewidth=1pt,
    innerleftmargin=6pt,innerrightmargin=6pt,
    innertopmargin=6pt,innerbottommargin=6pt]
\textbf{\small AT A GLANCE}\par\vspace{4pt}
\begin{description}[leftmargin=0pt,labelsep=4pt,itemsep=2pt,parsep=0pt,
    font=\small\bfseries]
  \item[Problem:] Mamba's fixed-size recurrent state overwrites old context as
    sequences grow, creating a representation bottleneck that limits long-range recall.
  \item[Why it matters:] Long-document QA, many-shot ICL, and genomic sequence modeling
    require retaining information across scales far beyond the recurrent state capacity.
  \item[Method:] Hierarchical Memory Mamba (HMM) extracts paragraph-level semantics
    from the hidden state and stores them in a persistent long-term memory for
    task-relevant retrieval.
  \item[Result:] Large improvements on Long-Range Arena, SCROLLS benchmarks, and
    many-shot ICL tasks; near-linear complexity vs.\ quadratic for Transformers.
  \item[Evidence:] Evaluated on six Long-Range Arena tasks, four SCROLLS benchmarks,
    and PG-19 language modeling up to 32k tokens.
\end{description}
\end{mdframed}

\section{The Big Picture}

Recurrent sequence models like Mamba process tokens in $O(N)$ time and $O(1)$ memory
per step, making them far cheaper than Transformers for long sequences. But the constant
memory comes with a hard limit: the recurrent hidden state $h_t \in \mathbb{R}^d$ must
compress \emph{all} context seen so far into a fixed-size vector at every step. On
short sequences, this works well. On sequences of thousands or tens of thousands of
tokens, earlier context is progressively overwritten---a phenomenon the authors call
the \emph{representation bottleneck}.

The problem is architectural, not a matter of training. No amount of fine-tuning on
long sequences fixes the fundamental capacity limit of a fixed recurrent state.
HMM attacks the bottleneck directly by introducing a structured memory hierarchy
above the recurrent state, inspired by the cognitive science of human memory
(sensory $\to$ working $\to$ long-term memory).

\section{Why Existing Approaches Fall Short}

\textbf{Larger recurrent state.} Expanding $d$ reduces the bottleneck but increases
per-step compute quadratically and does not scale gracefully.

\textbf{Mamba-2 / HGRN2.} These architectural improvements to SSMs increase state
capacity and improve practical long-context performance, but the fixed-state bottleneck
remains---performance still degrades on Path-X and other tasks requiring recall beyond
the effective state horizon.

\textbf{Hybrid Mamba-Transformer models} (Jamba, Zamba) insert full attention layers
at regular intervals, recovering long-range recall at the cost of $O(N^2)$ for those
layers. HMM achieves near-linear complexity throughout by using sparse key-value
retrieval rather than full attention for the long-term memory.

\textbf{KV-cache SSMs.} Some recent models cache a rolling window of past hidden
states and attend over them (e.g.\ RWKV-7's chunk attention). Window-limited KV
caches still impose a horizon; HMM's long-term memory has no fixed horizon.

\section{Core Mechanism}

HMM adds two memory tiers to a standard Mamba backbone:

\textbf{Tier 1---Sensory Memory (Mamba hidden state).}
The standard recurrent state $h_t$ processes the current token in the context of
recent tokens. It is fast and fine-grained but has fixed capacity $d$.

\textbf{Tier 2---Working Memory (PLS extractor).}
Every $P$ tokens (a ``paragraph''), a lightweight cross-attention module reads the
$P$ hidden states $h_{[kP:(k+1)P]}$ and compresses them into a single
\emph{Paragraph-Level Semantic} (PLS) vector:
\[
\mathrm{PLS}_k = \mathrm{MultiHead}\!\left(
  Q_{\mathrm{probe}},\;
  K = W_K H_k,\;
  V = W_V H_k
\right),
\]
where $H_k = [h_{kP}, \ldots, h_{(k+1)P-1}] \in \mathbb{R}^{P \times d}$ and
$Q_{\mathrm{probe}} \in \mathbb{R}^{q \times d}$ is a learned query matrix
($q \ll P$). The PLS captures the slow, paragraph-scale semantics that the fast
sensory memory cannot retain.

\textbf{Tier 3---Long-Term Memory (LTM store).}
PLS vectors accumulate in a growing bank $\mathcal{M} = \{(\mathrm{PLS}_k, k)\}$.
At each step, the current hidden state queries the bank via $k$-NN retrieval
followed by a small attention:
\[
\mathrm{retrieved}_t = \mathrm{Attn}\!\left(
  W_q h_t,\; K_{\mathrm{LTM}},\; V_{\mathrm{LTM}}
\right),
\]
and the retrieved content is injected into the Mamba state update via a gating
mechanism:
\[
h'_t = h_t + \sigma(W_g h_t) \odot \mathrm{retrieved}_t.
\]

\section{Technical Formulation}

The complexity of HMM for a sequence of $N$ tokens with paragraph size $P$:

\textbf{Sensory processing:} $O(Nd)$ (linear, standard Mamba).

\textbf{PLS extraction:} One cross-attention over $P$ tokens per paragraph:
$O((N/P) \cdot P \cdot q \cdot d) = O(Nqd/1)$, effectively $O(N)$ for fixed $q$.

\textbf{LTM retrieval:} $N/P$ entries in the bank; each retrieval is $O(N/P)$ for
exact attention over all entries, giving $O(N \cdot N/P) = O(N^2/P)$ total.
For $P \gg 1$ (e.g.\ $P = 256$), this is orders of magnitude cheaper than full
attention ($O(N^2)$).

The training objective adds no new loss terms: HMM is trained end-to-end on the
next-token prediction objective, with the PLS extractor and LTM read/write
operations differentiable throughout.

\section{Learning or Inference Procedure}

\textbf{Initialization.} HMM can be initialized from a pretrained Mamba checkpoint;
only the PLS extractor and LTM projection matrices are new parameters (typically
$<$5\% of total parameters).

\textbf{Fine-tuning.} Continued pre-training on long-document corpora for 10--20k
steps suffices to activate the LTM pathway; full long-context fine-tuning with HMM
from scratch is also supported.

\textbf{Inference.} The LTM bank grows as the sequence is processed; retrieval uses
approximate nearest neighbors (FAISS) for bank sizes above 10k entries. The $k$-NN
search cost is sub-linear in practice due to FAISS's IVF indexing.

\section{What the Guarantee Says}

\textbf{Monotone recall.} A PLS vector is written for every paragraph in the input
without overwriting; the LTM bank is append-only. This guarantees that information
from paragraph $k$ is always retrievable from step $k' > k$, unlike the sensory
state which is overwritten at step $k+1$.

\textbf{Complexity bound.} For paragraph size $P = O(\sqrt{N})$, LTM retrieval cost
is $O(N^{3/2})$---between linear and quadratic. For fixed $P$, it is $O(N^2/P)$,
which for $P = 256$ and $N = 32{,}768$ is $128\times$ cheaper than full attention.

\section{Experimental Findings}

\begin{itemize}[itemsep=2pt,parsep=0pt]
  \item \textbf{Long-Range Arena:} HMM outperforms Mamba-2 by 4--9 points on Path-X
    and Retrieval tasks, where long-range recall is decisive.
  \item \textbf{SCROLLS (QuALITY):} 6-point accuracy improvement over Mamba-2 on
    multi-hop reading comprehension requiring cross-paragraph reasoning.
  \item \textbf{Many-shot ICL:} At 64-shot, HMM accuracy is 3 points above Mamba-2
    and continues improving; Mamba-2 saturates and slightly degrades.
  \item \textbf{PG-19 language modeling:} 0.4 perplexity reduction at 16k tokens;
    0.9 reduction at 32k tokens, showing increasing benefit at longer contexts.
  \item \textbf{Parameter overhead:} $<$5\% additional parameters for the PLS
    extractor and LTM projections.
\end{itemize}

\section{Ablations and Interpretation}

\textbf{Paragraph size $P$.} Smaller $P$ gives richer PLS resolution but increases
LTM bank growth rate and retrieval cost. $P = 128$--$256$ gives the best quality/cost
tradeoff in experiments.

\textbf{LTM vs.\ no LTM.} Removing the LTM (keeping only PLS extraction) gives
intermediate improvement over base Mamba but still degrades on the longest sequences,
confirming that persistent retrieval---not just better compression---is needed.

\textbf{Retrieval head count.} Using 4 retrieval attention heads vs.\ 1 improves
SCROLLS accuracy by 2 points; beyond 4 heads, gains are marginal.

\section*{Reference}

Qinwen Wang, Jieping Luo, Aoxiang Qin, Ruoyu Zhao, Jianxiong Tang, Wei Zhang,
Zhichao Lu, Luziwei Leng.
\textbf{Mamba with Hierarchical Memory: Solving Representation Bottleneck in Long
Sequence Modeling}. arXiv:2608.02347, August 2026.
\url{https://arxiv.org/abs/2608.02347}

\end{document}
